\chapter{Quantum Computing Foundations}
\label{ch:quantum}

This chapter assembles the quantum-computing concepts on which
Chapter~\ref{ch:qml} builds: how quantum information is represented, how it is
transformed and measured, what today's hardware can realistically execute, and
where classical simulation of quantum circuits reaches its limits. The
treatment is deliberately focused; comprehensive accounts can be found in
\citet{nielsen2010quantum}.

\section{Qubits, superposition, and entanglement}
\label{sec:qc-qubits}

The elementary unit of quantum information is the \emph{qubit}, a two-level
quantum system. Its state is a unit vector in the two-dimensional complex
Hilbert space $\Hil_1 = \C^2$, written in Dirac notation as
\begin{equation}
  \ket{\psi} = \alpha \ket{0} + \beta \ket{1},
  \qquad \alpha, \beta \in \C, \quad
  \abs{\alpha}^2 + \abs{\beta}^2 = 1,
  \label{eq:qubit}
\end{equation}
where $\{\ket{0}, \ket{1}\}$ is the \emph{computational basis}. A state with
both $\alpha \neq 0$ and $\beta \neq 0$ is a \emph{superposition}: not a
probabilistic mixture of $\ket{0}$ and $\ket{1}$, but a distinct state whose
complex amplitudes can interfere. Global phase is unobservable, so a single
qubit state is fully described by two real angles,
\begin{equation}
  \ket{\psi} = \cos\tfrac{\vartheta}{2}\,\ket{0}
             + e^{i\varphi} \sin\tfrac{\vartheta}{2}\,\ket{1},
  \qquad \vartheta \in [0, \pi],\ \varphi \in [0, 2\pi),
\end{equation}
which identifies the state with a point on the \emph{Bloch sphere}. This
picture matters for machine learning because the most common way to load a
real-valued feature into a qubit --- \emph{angle encoding},
Section~\ref{sec:qml-encoding} --- is precisely to use the feature as a Bloch
rotation angle; the $[0,\pi]$ feature scaling adopted throughout this thesis
maps each feature onto a great arc of the sphere without wrapping.

A register of $n$ qubits lives in the tensor-product space
$\Hil_n = (\C^2)^{\otimes n} \cong \C^{2^n}$: a general state
\begin{equation}
  \ket{\psi} = \sum_{z \in \{0,1\}^n} c_z \ket{z},
  \qquad \sum_z \abs{c_z}^2 = 1
  \label{eq:nqubit}
\end{equation}
carries $2^n$ complex amplitudes. This exponential dimensionality is the
resource QML hopes to exploit --- and, simultaneously, the reason
$n \lesssim 25$-qubit systems remain classically simulable
(Section~\ref{sec:qc-simulability}).

A multi-qubit state that cannot be written as a tensor product of single-qubit
states is \emph{entangled}; the canonical example is the Bell state
$\ket{\Phi^+} = (\ket{00} + \ket{11})/\sqrt{2}$. Entanglement is generated by
two-qubit gates and is the structural ingredient that distinguishes genuinely
quantum feature maps and ansätze from products of independent single-qubit
operations. Whether that ingredient actually helps on standard learning tasks
is an empirical question this thesis probes directly with
entanglement-on/off ablations (Chapter~\ref{ch:experiments}), motivated by
benchmark findings that removing entanglement often costs little
\citep{bowles2024better}.

\section{Gates and circuits}
\label{sec:qc-gates}

Closed-system quantum evolution is unitary: a \emph{gate} on $k$ qubits is a
unitary $U \in \mathrm{U}(2^k)$ acting on the corresponding factors of
$\Hil_n$. The single-qubit gates used constantly in this thesis are the Pauli
operators
\begin{equation}
  X = \begin{pmatrix} 0 & 1 \\ 1 & 0 \end{pmatrix}, \quad
  Y = \begin{pmatrix} 0 & -i \\ i & 0 \end{pmatrix}, \quad
  Z = \begin{pmatrix} 1 & 0 \\ 0 & -1 \end{pmatrix},
\end{equation}
the Hadamard gate $H = \tfrac{1}{\sqrt{2}}\bigl(\begin{smallmatrix} 1 & 1 \\
1 & -1 \end{smallmatrix}\bigr)$, which maps computational-basis states to
uniform superpositions, and the parameterized rotations
\begin{equation}
  R_P(\theta) = e^{-i \theta P / 2}
              = \cos\tfrac{\theta}{2}\, I - i \sin\tfrac{\theta}{2}\, P,
  \qquad P \in \{X, Y, Z\}.
  \label{eq:rotation}
\end{equation}
Rotations are the workhorses of variational circuits: data values enter as
rotation angles (encoding), and trainable parameters $\bm{\theta}$ enter as
rotation angles (ansatz). The form \eqref{eq:rotation}, with generator
satisfying $P^2 = I$, is also what makes the parameter-shift rule of
Section~\ref{sec:qml-vqa} exact.

Entangling operations are typically the controlled-NOT
($\mathrm{CNOT}\,\ket{a,b} = \ket{a, a \oplus b}$) or controlled-$Z$ gates.
Together with single-qubit rotations they form a universal set: any unitary on
$n$ qubits can be approximated to arbitrary precision by a finite circuit over
this set \citep{nielsen2010quantum}.

A \emph{quantum circuit} is a sequence of gates applied to an initial state,
conventionally $\ket{0}^{\otimes n}$. Two cost measures matter throughout this
thesis. The \emph{gate count} determines simulation effort and hardware error
accumulation; the \emph{depth} --- the number of layers of gates that cannot be
parallelized --- determines how long a computation must survive decoherence.
On hardware, circuits must additionally be \emph{transpiled}: rewritten over
the device's native gate set and routed to respect its qubit-connectivity
graph, a step that can substantially inflate both counts and is logged
explicitly in the resource-accounting experiments (E9).

\section{Measurement and expectation values}
\label{sec:qc-measurement}

Information leaves a quantum computer only through measurement. Measuring all
qubits of state \eqref{eq:nqubit} in the computational basis returns bitstring
$z$ with probability $p(z) = \abs{c_z}^2$ (the \emph{Born rule}), and the state
collapses accordingly; a measurement is not a passive read-out but a
destructive sampling operation.

Machine-learning pipelines almost always consume \emph{expectation values} of
observables rather than raw bitstrings. For a Hermitian observable $M$ and
state $\ket{\psi}$, the quantity of interest is
\begin{equation}
  \langle M \rangle = \bra{\psi} M \ket{\psi},
\end{equation}
for example $M = Z_1$ (the Pauli-$Z$ on a designated readout qubit), whose
expectation in $[-1, 1]$ serves directly as a binary-classification score.
On real devices and shot-based simulators, $\langle M \rangle$ is estimated by
repeating the circuit $S$ times (``shots'') and averaging outcomes. The
estimator's standard error scales as
\begin{equation}
  \varepsilon \;\sim\; \sqrt{\frac{\Var[M]}{S}} \;=\; O\!\bigl(S^{-1/2}\bigr),
  \label{eq:shotnoise}
\end{equation}
so each additional decimal digit of precision costs a hundredfold more shots.
Equation~\eqref{eq:shotnoise} is the root of two phenomena central to this
thesis: the \emph{shot-noise regime} of experiment E8, where model quality
degrades gracefully as $S$ shrinks, and the interaction of sampling error with
the \emph{concentration} effects of Section~\ref{sec:qml-trainability}, where
signals that decay exponentially in the qubit count fall below any affordable
shot budget.

\section{The NISQ era: noise and error mitigation}
\label{sec:qc-nisq}

Current devices --- tens to a few hundred qubits, no fault tolerance ---
define the \emph{noisy intermediate-scale quantum} (NISQ) era
\citep{preskill2018nisq}. Physical qubits decohere: excited states relax
toward $\ket{0}$ on a timescale $T_1$, and phase coherence between $\ket{0}$
and $\ket{1}$ decays on a timescale $T_2$. Gates are imperfect, with two-qubit
gates typically an order of magnitude noisier than single-qubit ones, and
readout itself misassigns outcomes with percent-level probability.

Noise is modeled by quantum channels acting on density matrices
$\rho = \ketbra{\psi}{\psi}$ (and their mixtures). The depolarizing channel
with error probability $p$,
\begin{equation}
  \mathcal{E}_{\mathrm{dep}}(\rho)
    = (1 - p)\,\rho + \frac{p}{3}\left( X\rho X + Y\rho Y + Z\rho Z \right),
\end{equation}
is the generic ``white noise'' model; bit-flip, phase-flip, and
amplitude-damping channels capture more specific error mechanisms, and
calibrated device models compose gate-level channels with measured $T_1$,
$T_2$, gate-error, and readout-error parameters of a physical backend. The
noise ladder of experiment E8 uses exactly this hierarchy: exact statevector
$\rightarrow$ finite shots $\rightarrow$ calibrated device-noise simulation
$\rightarrow$ real hardware.

Full quantum error \emph{correction} is beyond NISQ budgets, but lightweight
error \emph{mitigation} is not. This thesis uses measurement-error (readout)
mitigation --- characterizing the readout-confusion matrix and applying its
(regularized) inverse to measured distributions --- for the hardware runs, and
notes zero-noise extrapolation as a standard alternative; heavier schemes are
out of scope. Mitigation reduces bias in expectation values at the price of
increased variance, i.e., more shots \citep{preskill2018nisq}.

\section{Limits of classical simulability}
\label{sec:qc-simulability}

Storing the state vector \eqref{eq:nqubit} of $n$ qubits requires $2^n$
complex amplitudes: at double precision, roughly $16 \times 2^n$ bytes ---
about 4\,GB at $n = 28$ and about 1\,TB at $n = 36$. Statevector simulation of
the circuit sizes in this thesis ($n \le 20$) is therefore comfortable on a
laptop, which is precisely why the experimental programme can run locally,
exactly, and reproducibly. Restricted circuit families admit far larger
simulations --- Clifford circuits via the stabilizer formalism, and
low-entanglement circuits via tensor-network contraction
\citep{nielsen2010quantum} --- and modern classical surrogate and
dequantization results (Section~\ref{sec:related-surrogates}) push the
simulable frontier further for many structured QML models specifically.

Two consequences frame everything that follows. First, the honest-framing
position of Section~\ref{sec:intro-framing} is not a caveat but a mathematical
fact: at the scales studied here, quantum models \emph{are} classical
computations, executed through an unusual parameterization. Second,
simulability is an asset for methodology: exact simulation provides
noise-free ground truth against which every shot, noise, and hardware effect
in E8 can be measured --- a separation of concerns impossible in
hardware-first studies.
